% !TEX root = ../../main.tex

\section{Map of Core Technical Terms}
\label{app:jargon-map}

This appendix maps the main technical terms used in the manuscript to their
plain-language meanings, their first point of entry in the paper, and their
role in the method. The goal is not to replace the formal development in the
main text, but to make the paper easier to navigate for readers who do not work
routinely with hierarchical clustering or projected test statistics.

\paragraph{Hierarchical clustering.}
This term first appears in the opening paragraph of the introduction. In the
manuscript, it means a family of methods that builds nested groups by
repeatedly merging similar samples or sample groups. Its role is to provide the
initial hierarchy from which candidate cluster boundaries are proposed.

\paragraph{Tree.}
This term also first appears in the opening paragraph of the introduction. Here
it means a rooted diagram of nested sample groups. Its role is to provide the
central geometric object on which all later statistical decisions are made.

\paragraph{Internal node.}
This term first appears in the opening paragraph of the introduction. It means
a merged group of descendant samples rather than a new observation. Its role is
to mark a candidate decision point at which the method may keep or reject a
split.

\paragraph{Root.}
This term first appears in the opening paragraph of the introduction. It means
the top node containing all samples. Its role is to anchor the hierarchy and
the top-down traversal that forms the final clusters.

\paragraph{Split.}
This term first appears in the second paragraph of the introduction. It means a
proposed division of one parent group into two child groups. Its role is to
define the local decision problem tested at each internal node.

\paragraph{Linkage tree.}
This term first appears in the second paragraph of the introduction. It means
the hierarchy returned by the agglomerative merging rule itself. Its role is to
distinguish the raw clustering proposal from the later statistical filtering
stage.

\paragraph{Subtree.}
This term first appears in the second paragraph of the introduction. It means a
node together with all descendant leaves below it. Its role is to define the
group object whose empirical feature distribution is tested.

\paragraph{Child-parent edge test.}
This term first appears in the third paragraph of the introduction. It means a
test asking whether a child subtree differs from the broader background
represented by its parent. Its role is to supply the first condition for
keeping a split.

\paragraph{Sibling test.}
This term first appears in the third paragraph of the introduction. It means a
test asking whether the two child subtrees under the same parent differ from
each other. Its role is to supply the second condition for keeping a split.

\paragraph{Multiple-testing correction.}
This term first appears in the third paragraph of the introduction. It means an
adjustment for the fact that many candidate splits are tested across the same
hierarchy. Its role is to control false discoveries created by repeated testing
across the tree.

\paragraph{Nested variance.}
This term first appears in the final paragraph of the introduction. It means a
variance formula for a comparison in which the child group is contained inside
the parent group. Its role is to make the child-parent edge test respect the
dependence created by subtree containment.

\paragraph{Parent-local summary of the feature space.}
This phrase first appears in the final paragraph of the introduction. It means
a lower-dimensional description built from the variation seen below one parent
node. Its role is to prevent the edge statistic from overcounting correlated
high-dimensional directions.

\paragraph{Post-selection calibration.}
This phrase first appears in the final paragraph of the introduction. It means
a correction for bias created when the same data both choose a split and test
it. Its role is to motivate the extra calibration layer in the sibling stage.

\paragraph{Traversal.}
This term first appears in the final paragraph of the introduction. It means a
top-down walk through the tree that decides where splitting stops. Its role is
to convert nodewise evidence into the final partition of samples.

\paragraph{Post-selection inflation.}
This phrase first appears in the method overview subsection. It means upward
bias in a test statistic caused by evaluating a split that was chosen from the
same data. Its role is to explain why raw sibling statistics are
anti-conservative and require deflation.

\paragraph{PCA directions.}
This phrase first appears in the sibling-test subsection on the projected
sibling statistic. It means principal-component directions summarizing the
dominant local variation at a parent node. Its role is to provide the main
projection basis used to reduce the sibling contrast before testing.

\paragraph{Satterthwaite calibration.}
This phrase first appears in the sibling-test subsection on reference-law
calibration. It means a moment-matching approximation that replaces a
complicated null law by a scaled chi-square distribution. Its role is to
produce the reference distribution used for the projected sibling statistic.

\paragraph{Benjamini--Hochberg correction.}
This phrase first appears in the sibling-test subsection on calibration
weights. It means a false-discovery-rate adjustment applied to a family of
$p$-values. Its role is to define how edge-test significance information is
summarized before sibling calibration.

\paragraph{False-discovery rate.}
This phrase first appears in the sibling-test subsection on runtime correction.
It means the expected proportion of false positives among declared discoveries.
Its role is to describe the operating error criterion used for the final
sibling-stage multiplicity adjustment.
